\documentclass{article}
\begin{document}
\section{Introduction}
Flow boiling in microchannels removes large heat loads from small areas. However, flow instabilities limit its use. Pressure oscillations can trigger early dryout. In this work, we test whether inlet restrictors damp these oscillations.

\section{Experimental method}
We built a test section with ten parallel channels. Each channel is 0.5 mm wide and 1 mm deep. We heat the base with a thin-film heater and record wall temperatures with eight thermocouples. We use deionised water at atmospheric pressure.

\section{Results and discussion}
Figure 4 shows the pressure drop signal with and without restrictors. Without restrictors, the signal oscillates strongly. With restrictors, the amplitude drops sharply. We see the same trend at all heat fluxes we tested.

Next, we examine the heat transfer coefficient. Figure 5 shows it against vapour quality. The coefficient rises at low quality and then falls. Restrictors shift the peak to higher quality. We think this is because the flow distributes more evenly between channels.

We also recorded high-speed videos. They show reversed flow without restrictors. Vapour slugs move back into the inlet plenum. With restrictors, we rarely see reversed flow. The videos therefore support the pressure data.

We now discuss the cost of this approach. Restrictors raise the total pressure drop. The pump must work harder. However, the gain in stability outweighs this cost for the conditions we tested.

\section{Conclusions}
In this work, we showed that inlet restrictors damp flow oscillations in parallel microchannels. They delay dryout and shift the peak heat transfer coefficient. The extra pressure drop is modest for the conditions we tested.
\end{document}
